
\documentclass{article}
\usepackage{colortbl}
\usepackage{makecell}
\usepackage{multirow}
\usepackage{supertabular}

\begin{document}

\newcounter{utterance}

\centering \large Interaction Transcript for game `cladder', experiment `full\_v1.5\_default', episode 855 with Qwen3.5{-}9B.
\vspace{24pt}

{ \footnotesize  \setcounter{utterance}{1}
\setlength{\tabcolsep}{0pt}
\begin{supertabular}{c@{$\;$}|p{.15\linewidth}@{}p{.15\linewidth}p{.15\linewidth}p{.15\linewidth}p{.15\linewidth}p{.15\linewidth}}
    \# & \multicolumn{2}{c}{Player} && \multicolumn{2}{c}{Game Master} \\
    \hline

    \theutterance \stepcounter{utterance}  
    & & & \multicolumn{4}{p{0.6\linewidth}}{
        \cellcolor[rgb]{0.9,0.9,0.9}{
            \makecell[{{p{\linewidth}}}]{
                \texttt{\tiny{[P1$\langle$GM]}}
                \texttt{You are an expert in causal inference. The following question is not a typical commonsense query, but rather a meticulously designed question created by a professor specializing in causal inference, intended to assess the students' mastery of the course content.} \\
\texttt{Imagine a self{-}contained, hypothetical world with only the following conditions, and without any unmentioned factors or causal relationships: Personality has a direct effect on relationship status. Appearance has a direct effect on relationship status. The overall probability of kindness is 85\%. For people with mean personalities and single people, the probability of attractive appearance is 38\%. For people with mean personalities and in a relationship, the probability of attractive appearance is 71\%. For people with kind personalities and single people, the probability of attractive appearance is 32\%. For people with kind personalities and in a relationship, the probability of attractive appearance is 57\%. If we look at people in a relationship, does the chance of attractive appearance decrease when kindness?} \\
\texttt{Start your answer with 'Yes' or 'No', followed by additional reasoning or evidence to support your explanation.} \\
            }
        }
    }
     \\ \\

    \theutterance \stepcounter{utterance}  
    & \multicolumn{4}{p{0.6\linewidth}}{
        \cellcolor[rgb]{0.561,0.737,0.561}{
            \makecell[{{p{\linewidth}}}]{
                \texttt{\tiny{[P1$\rangle$GM]}}
                \texttt{No.} \\
\\ 
\texttt{**Reasoning:**} \\
\\ 
\texttt{To determine whether the chance of attractive appearance decreases when kindness increases among people in a relationship, we must calculate the conditional probability of having an attractive appearance given that a person is in a relationship, for both "kind" and "mean" personality types.} \\
\\ 
\texttt{Let \$A\$ be the event of having an attractive appearance, \$R\$ be the event of being in a relationship, \$K\$ be the event of having a kind personality, and \$M\$ be the event of having a mean personality.} \\
\\ 
\texttt{We are given the following probabilities:} \\
\texttt{1.  \$P(A | M, R) = 71\textbackslash{}\%\$ (Probability of attractive appearance for mean personalities in a relationship).} \\
\texttt{2.  \$P(A | K, R) = 57\textbackslash{}\%\$ (Probability of attractive appearance for kind personalities in a relationship).} \\
\\ 
\texttt{The question asks: "If we look at people in a relationship, does the chance of attractive appearance decrease when kindness?"} \\
\\ 
\texttt{Comparing the two conditional probabilities within the subset of people in a relationship (\$R\$):} \\
\texttt{*   For mean personalities (\$M\$), the probability is **71\%**.} \\
\texttt{*   For kind personalities (\$K\$), the probability is **57\%**.} \\
\\ 
\texttt{Since \$57\textbackslash{}\% < 71\textbackslash{}\%\$, the probability of having an attractive appearance is lower for people with kind personalities compared to those with mean personalities, specifically within the group of people} \\
            }
        }
    }
    & & \\ \\

    \theutterance \stepcounter{utterance}  
    & & & \multicolumn{2}{p{0.3\linewidth}}{
        \cellcolor[rgb]{0.9,0.9,0.9}{
            \makecell[{{p{\linewidth}}}]{
                \texttt{\tiny{[GM$|$GM]}}
                \texttt{No} \\
            }
        }
    }
    & & \\ \\

    \theutterance \stepcounter{utterance}  
    & & & \multicolumn{2}{p{0.3\linewidth}}{
        \cellcolor[rgb]{0.9,0.9,0.9}{
            \makecell[{{p{\linewidth}}}]{
                \texttt{\tiny{[GM$|$GM]}}
                \texttt{yes} \\
            }
        }
    }
    & & \\ \\

    \theutterance \stepcounter{utterance}  
    & & & \multicolumn{2}{p{0.3\linewidth}}{
        \cellcolor[rgb]{0.9,0.9,0.9}{
            \makecell[{{p{\linewidth}}}]{
                \texttt{\tiny{[GM$|$GM]}}
                \texttt{game\_result = LOSE} \\
            }
        }
    }
    & & \\ \\

\end{supertabular}
}

\end{document}
